\section{Fixed packages saturate; growth requires switching}\label{sec:fixed-package}

This section makes precise the intuition that a fixed closure rule cannot by itself keep producing new material. The results are elementary, but the exact hypotheses are informative, and the quantitative bound of Theorem~\ref{thm:switch-budget} is sharp.

\subsection{Saturation}

\begin{proposition}[One-step saturation]\label{prop:one-step-saturation}
Let $p$ be a package on $X$. For every $x\in X$ and every $n\in\N$, $p^{[n+1]}(x)=p(x)$.
\end{proposition}

\begin{proof}
Induction on $n$. For $n=0$ both sides are $p(x)$. If $p^{[n+1]}(x)=p(x)$, then $p^{[n+2]}(x)=p\bigl(p^{[n+1]}(x)\bigr)=p(p(x))=p(x)$ by idempotence.
\leanref{\lid{closure\_iterate\_saturates\_one\_step}}
\end{proof}

Idempotence is exactly what rules out further growth. The following characterization makes that precise.

\begin{theorem}[Idempotence is absence of growth]\label{thm:idempotent-iff}
A map $p\colon X\to X$ is idempotent if and only if no point of $X$ is a persistent growth witness for~$p$.
\end{theorem}

\begin{proof}
If $p$ is idempotent and $n\ge 1$, Proposition~\ref{prop:one-step-saturation} gives $p^{[n+1]}(x)=p(x)=p^{[n]}(x)$. Conversely, if $p(p(x))\neq p(x)$ for some $x$, then $x$ is a witness with $n=1$.
\leanref{\lid{frozen\_package\_no\_persistent\_growth}, \lid{idempotent\_iff\_no\_persistent\_growth}}
\end{proof}

The hypothesis cannot be weakened to ``the map stays fixed''. The successor map on $\N$ is fixed, yet every point is a growth witness. Growth from a fixed rule is therefore not a contradiction in general. It is ruled out for packages because they are idempotent.

\subsection{Package change, with a witness}

If growth does occur, something must have changed. The next result locates the change.

\begin{theorem}[Witnessed package change]\label{thm:witnessed-change}
Let $p,q$ be packages on $X$ and $x\in X$. If $q(p(x))\neq p(x)$, then $p$ and $q$ differ, and $z=p(x)$ is a change witness: $p(z)\neq q(z)$. Moreover, for any packages $p,q$, we have $p\neq q$ if and only if a change witness exists.
\end{theorem}

\begin{proof}
By idempotence $p(z)=p(p(x))=p(x)=z$, while $q(z)\neq z$ by hypothesis. Hence $p(z)\neq q(z)$, and in particular $p\neq q$. The equivalence is function extensionality: two maps are equal exactly when they agree at every point.
\leanref{\lid{persistent\_growth\_witness\_implies\_package\_change}, \lid{package\_changed\_iff\_witness}}
\end{proof}

Theorem~\ref{thm:witnessed-change} identifies \emph{where} two packages differ, not \emph{why}. In Six Birds terms, package change is where the rewrite and gating primitives (P1, P2) are read as acting. That reading is an interpretation, and nothing in the theorem singles out one mechanism.

\subsection{A quantitative switch budget}

Now follow a whole schedule (Definition~\ref{def:schedule}). The first step may change the state for free, because a package can move a point that is not yet closed. After that, each state change must be paid for by a switch at the preceding index.

\begin{lemma}[Each later change forces a switch]\label{lem:orbit-switch}
For every schedule, initial state and $t\in\N$: if $x_{t+2}\neq x_{t+1}$, then $p_t\neq p_{t+1}$, with change witness $x_{t+1}$.
\end{lemma}

\begin{proof}
Apply Theorem~\ref{thm:witnessed-change} with $p=p_t$, $q=p_{t+1}$ and $x=x_t$. Then $p(x)=x_{t+1}$ and $q(p(x))=x_{t+2}$.
\leanref{\lid{orbit\_change\_requires\_switch}}
\end{proof}

\begin{theorem}[Switch budget]\label{thm:switch-budget}
For every schedule of packages, every initial state and every horizon $n\in\N$,
\[
\mathrm{changes}_n\;\le\;1+\mathrm{switches}_n .
\]
\end{theorem}

\begin{proof}
Let $S$ be the set of change steps $t\le n$ with $t\ge 1$. By Lemma~\ref{lem:orbit-switch}, $t\mapsto t-1$ maps $S$ into the set of switch indices below~$n$. This map is injective. Step $0$ contributes at most one further change.
\leanref{\lid{package\_switch\_budget}}
\end{proof}

\begin{corollary}[Unbounded growth needs unbounded switching]\label{cor:unbounded}
If an orbit changes state at arbitrarily late steps, then the schedule switches packages at arbitrarily late indices.
\leanref{\lid{unbounded\_changes\_require\_unbounded\_switches}}
\end{corollary}

The corollary counts \emph{switches}, not \emph{distinct packages}. The next result shows that two packages are enough for unbounded growth, and that the budget cannot be improved.

\begin{proposition}[Sharpness]\label{prop:sharpness}
Define $o,\,e\colon\N\to\N$ by
\[
o(n)=2\lfloor n/2\rfloor+1,\qquad e(n)=2\lfloor (n+1)/2\rfloor ,
\]
which round up to the nearest odd and the nearest even number, respectively. Both maps are computable, monotone, inflationary ($n\le o(n)$, $n\le e(n)$) and idempotent. For the alternating schedule $p_t=o$ when $t$ is even and $p_t=e$ when $t$ is odd, the orbit from $x_0=0$ is $x_t=t$ for every $t$. Consequently $\mathrm{changes}_n=n+1$ and $\mathrm{switches}_n=n$, so the bound of Theorem~\ref{thm:switch-budget} holds with equality at every horizon.
\end{proposition}

\begin{proof}
The four laws are routine arithmetic on $\lfloor\cdot/2\rfloor$, and both maps are primitive recursive. If $x_t=t$ with $t$ even, then $x_{t+1}=o(t)=t+1$. If $t$ is odd, then $x_{t+1}=e(t)=t+1$. Every step changes the state, and $o\neq e$ (for instance $o(0)=1\neq 0=e(0)$), so every index is a switch.
\leanref{\lid{odd\_round\_laws}, \lid{even\_round\_laws}, \lid{odd\_round\_computable}, \lid{even\_round\_computable}, \lid{alternating\_orbit\_exact}, \lid{alternating\_budget\_sharp}}
\end{proof}

\begin{figure}[t]
\centering
% Figure 2: orbits of the odd/even rounding closures (Proposition "Sharpness")
\begin{tikzpicture}[x=0.95cm, y=0.48cm, font=\footnotesize]
  % grid and axes
  \foreach \y in {0,2,4,6,8} \draw[black!12, line width=0.4pt] (0,\y) -- (8.4,\y);
  \draw[-{Stealth[length=2mm]}, line width=0.6pt] (0,0) -- (8.6,0) node[right] {$t$};
  \draw[-{Stealth[length=2mm]}, line width=0.6pt] (0,0) -- (0,8.9) node[above] {$x_t$};
  \foreach \t in {0,...,8} \draw (\t,0) -- (\t,-0.15) node[below] {\t};
  \foreach \y in {2,4,6,8} \draw (0,\y) -- (-0.08,\y) node[left] {\y};
  % alternating schedule o,e,o,e,...: x_t = t
  \draw[seriesB, line width=1.1pt] (0,0) \foreach \t in {1,...,8} { -- (\t,\t) };
  \foreach \t in {0,...,8} \node[seriesB, fill=seriesB, rectangle, inner sep=1.7pt] at (\t,\t) {};
  % fixed schedule o,o,o,...: 0,1,1,1,...
  \draw[seriesA, line width=1.1pt] (0,0) -- (1,1) -- (8,1);
  \foreach \t/\y in {0/0,1/1,2/1,3/1,4/1,5/1,6/1,7/1,8/1}
     \node[seriesA, fill=white, draw=seriesA, circle, line width=0.9pt, inner sep=1.4pt] at (\t,\y) {};
  % switch ticks for the alternating schedule
  \foreach \t in {0,...,7} \draw[seriesB, line width=0.8pt] (\t+0.5,-1.45) -- (\t+0.5,-1.05);
  \node[seriesB, anchor=east, font=\scriptsize] at (-0.1,-1.25) {switches};
  % direct labels
  \node[seriesB, anchor=west, align=left] at (0.4,6.6) {alternate $o,e,o,e,\dots$:\\ every step changes};
  \node[seriesA, anchor=west] at (3.0,1.9) {fixed $o,o,o,\dots$: one move, then saturated};
\end{tikzpicture}

\caption{Orbits from $x_0=0$ of the two closures of Proposition~\ref{prop:sharpness}. Applying the odd-rounding closure $o$ at every step stops after one move (blue circles). Alternating $o$ and the even-rounding closure $e$ climbs one unit per step forever (orange squares). In the alternating run every index is a switch (tick marks), and the switch budget of Theorem~\ref{thm:switch-budget} is attained at every horizon.}
\label{fig:orbits}
\end{figure}

Figure~\ref{fig:orbits} contrasts the two behaviors. The example is deliberately simple. Its point is that unbounded growth needs no infinite supply of new rules, only infinitely many changes of rule.

\begin{example}[The consistency progression]\label{ex:consistency-progression}
This example is an informal illustration; it is not part of the Lean development. Let $T_0$ be a sound, effectively axiomatized extension of $\EA$ and put $T_{k+1}=T_k+\Con(T_k)$. Let $p_k$ be deductive closure under $T_k$, acting on sets of sentences: $p_k(S)$ is the set of consequences of $T_k\cup S$. Each $p_k$ is a closure operator, hence a package. Starting from $x_0=\varnothing$, the orbit is $x_{k+1}=$ the set of theorems of $T_k$, because $T_{k-1}\subseteq T_k$. Each $T_k$ is sound, hence consistent. By the second incompleteness theorem, $\Con(T_{k-1})$ lies in $x_{k+1}$ but not in $x_k$ for every $k\ge 1$, so every step changes the state. Lemma~\ref{lem:orbit-switch} then says that each $p_k$ differs from $p_{k+1}$ at the current theory. This is the package-level form of the familiar fact that no fixed deductive closure reaches its own consistency statement.
\end{example}
